\section{Inline Visual Route: Reader Routes}

These figures are placed in the main argument as visual proof-of-reading aids: each one separates objects, direction, quantitative or formal anchor, and evidence route.

\subsection{Reviewer attack route}

This figure is included because the reader must see how A promoted OC claim is constrained by The exact proof, finite witness, replay row, or comparator that caps it. The highlighted review variable is attack point count, and the formal anchor is $C_{\mathrm{public}}\Rightarrow E_{\mathrm{bounded}}$. The nearby prose names the proof, evidence row, table, replay check, or appendix that carries the claim.

\begin{figure}[!htbp]
\centering
\resizebox{0.96\textwidth}{!}{%
\begin{tikzpicture}[x=1cm,y=1cm,>=Latex,every node/.style={font=\small}]
  \definecolor{ocBlue}{HTML}{173B57}
\definecolor{ocGold}{HTML}{A77D2A}
\definecolor{ocGreen}{HTML}{4B7F52}
\definecolor{ocRed}{HTML}{8E3B32}
\definecolor{ocGray}{HTML}{ECEFF1}
  \draw[rounded corners=10pt,fill=ocGray!35,draw=ocBlue,line width=0.9pt] (-6.8,-3.4) rectangle (6.8,3.4);
  \node[font=\bfseries\large,ocBlue] at (0,3.0) {Reviewer attack route};
  \node[draw,rounded corners=5pt,fill=blue!7,text width=0.24\textwidth,align=center,minimum height=1.05cm] (a) at (-4.35,1.0) {A promoted OC claim};
  \node[draw,rounded corners=5pt,fill=green!8,text width=0.24\textwidth,align=center,minimum height=1.05cm] (b) at (4.35,1.0) {The exact proof, finite witness, replay row, or comparator that caps it};
  \draw[->,line width=1.0pt,ocGreen] (a.east) -- node[above,fill=ocGray!35,inner sep=2pt,align=center] {claim-to-anchor inspection} (b.west);
  \node[draw,rounded corners=5pt,fill=orange!10,text width=0.28\textwidth,align=center,minimum height=0.9cm] (r) at (-4.35,-1.45) {attack point count};
  \node[draw,rounded corners=5pt,fill=white,text width=0.28\textwidth,align=center,minimum height=0.9cm] (m) at (0,-1.45) {attack point count};
  \node[draw,rounded corners=5pt,fill=red!6,text width=0.28\textwidth,align=center,minimum height=0.9cm] (f) at (4.35,-1.45) {formal anchor: \( C_{\mathrm{public}}\Rightarrow E_{\mathrm{bounded}} \)};
  \draw[->,dashed,ocGold] (r.north) -- (a.south);
  \draw[->,dashed,ocGold] (m.north) -- ($(a)!0.5!(b)$);
  \draw[->,dashed,ocGold] (f.north) -- (b.south);
\end{tikzpicture}}
\caption{Reviewer attack route. This figure shows the reading relation from A promoted OC claim to The exact proof, finite witness, replay row, or comparator that caps it; the review variable is attack point count, the formal anchor is \( C_{\mathrm{public}}\Rightarrow E_{\mathrm{bounded}} \), and the evidence link is the named proof, table, replay row, or appendix section cited near the figure.}
\label{fig:r012-inline-28e_oc133_inline_figures_reader_routes-1}
\end{figure}

\subsection{Theorist model route}

This figure is included because the reader must see how Continuum intuition is constrained by Canonical tuple and K-level witness discipline. The highlighted review variable is model anchors, and the formal anchor is $K=(\Omega,\partial\Omega,A,\Theta,P,J,C,k,M)$. The nearby prose names the proof, evidence row, table, replay check, or appendix that carries the claim.

\begin{figure}[!htbp]
\centering
\resizebox{0.96\textwidth}{!}{%
\begin{tikzpicture}[x=1cm,y=1cm,>=Latex,every node/.style={font=\small}]
  \definecolor{ocBlue}{HTML}{173B57}
\definecolor{ocGold}{HTML}{A77D2A}
\definecolor{ocGreen}{HTML}{4B7F52}
\definecolor{ocRed}{HTML}{8E3B32}
\definecolor{ocGray}{HTML}{ECEFF1}
  \draw[rounded corners=10pt,fill=ocGray!35,draw=ocBlue,line width=0.9pt] (-6.8,-3.4) rectangle (6.8,3.4);
  \node[font=\bfseries\large,ocBlue] at (0,3.0) {Theorist model route};
  \node[draw,rounded corners=5pt,fill=blue!7,text width=0.24\textwidth,align=center,minimum height=1.05cm] (a) at (-4.35,1.0) {Continuum intuition};
  \node[draw,rounded corners=5pt,fill=green!8,text width=0.24\textwidth,align=center,minimum height=1.05cm] (b) at (4.35,1.0) {Canonical tuple and K-level witness discipline};
  \draw[->,line width=1.0pt,ocGreen] (a.east) -- node[above,fill=ocGray!35,inner sep=2pt,align=center] {concept-to-formal-object reading} (b.west);
  \node[draw,rounded corners=5pt,fill=orange!10,text width=0.28\textwidth,align=center,minimum height=0.9cm] (r) at (-4.35,-1.45) {model anchors};
  \node[draw,rounded corners=5pt,fill=white,text width=0.28\textwidth,align=center,minimum height=0.9cm] (m) at (0,-1.45) {model anchors};
  \node[draw,rounded corners=5pt,fill=red!6,text width=0.28\textwidth,align=center,minimum height=0.9cm] (f) at (4.35,-1.45) {formal anchor: \( K=(\Omega,\partial\Omega,A,\Theta,P,J,C,k,M) \)};
  \draw[->,dashed,ocGold] (r.north) -- (a.south);
  \draw[->,dashed,ocGold] (m.north) -- ($(a)!0.5!(b)$);
  \draw[->,dashed,ocGold] (f.north) -- (b.south);
\end{tikzpicture}}
\caption{Theorist model route. This figure shows the reading relation from Continuum intuition to Canonical tuple and K-level witness discipline; the review variable is model anchors, the formal anchor is \( K=(\Omega,\partial\Omega,A,\Theta,P,J,C,k,M) \), and the evidence link is the named proof, table, replay row, or appendix section cited near the figure.}
\label{fig:r012-inline-28e_oc133_inline_figures_reader_routes-2}
\end{figure}

\subsection{Applied architecture route}

This figure is included because the reader must see how Enterprise or AI system boundary is constrained by Domain-bounded projection and falsifier. The highlighted review variable is decision row, and the formal anchor is $P_D(K)\rightarrow F_D$. The nearby prose names the proof, evidence row, table, replay check, or appendix that carries the claim.

\begin{figure}[!htbp]
\centering
\resizebox{0.96\textwidth}{!}{%
\begin{tikzpicture}[x=1cm,y=1cm,>=Latex,every node/.style={font=\small}]
  \definecolor{ocBlue}{HTML}{173B57}
\definecolor{ocGold}{HTML}{A77D2A}
\definecolor{ocGreen}{HTML}{4B7F52}
\definecolor{ocRed}{HTML}{8E3B32}
\definecolor{ocGray}{HTML}{ECEFF1}
  \draw[rounded corners=10pt,fill=ocGray!35,draw=ocBlue,line width=0.9pt] (-6.8,-3.4) rectangle (6.8,3.4);
  \node[font=\bfseries\large,ocBlue] at (0,3.0) {Applied architecture route};
  \node[draw,rounded corners=5pt,fill=blue!7,text width=0.24\textwidth,align=center,minimum height=1.05cm] (a) at (-4.35,1.0) {Enterprise or AI system boundary};
  \node[draw,rounded corners=5pt,fill=green!8,text width=0.24\textwidth,align=center,minimum height=1.05cm] (b) at (4.35,1.0) {Domain-bounded projection and falsifier};
  \draw[->,line width=1.0pt,ocGreen] (a.east) -- node[above,fill=ocGray!35,inner sep=2pt,align=center] {application mapping} (b.west);
  \node[draw,rounded corners=5pt,fill=orange!10,text width=0.28\textwidth,align=center,minimum height=0.9cm] (r) at (-4.35,-1.45) {decision row};
  \node[draw,rounded corners=5pt,fill=white,text width=0.28\textwidth,align=center,minimum height=0.9cm] (m) at (0,-1.45) {decision row};
  \node[draw,rounded corners=5pt,fill=red!6,text width=0.28\textwidth,align=center,minimum height=0.9cm] (f) at (4.35,-1.45) {formal anchor: \( P_D(K)\rightarrow F_D \)};
  \draw[->,dashed,ocGold] (r.north) -- (a.south);
  \draw[->,dashed,ocGold] (m.north) -- ($(a)!0.5!(b)$);
  \draw[->,dashed,ocGold] (f.north) -- (b.south);
\end{tikzpicture}}
\caption{Applied architecture route. This figure shows the reading relation from Enterprise or AI system boundary to Domain-bounded projection and falsifier; the review variable is decision row, the formal anchor is \( P_D(K)\rightarrow F_D \), and the evidence link is the named proof, table, replay row, or appendix section cited near the figure.}
\label{fig:r012-inline-28e_oc133_inline_figures_reader_routes-3}
\end{figure}

\subsection{Strategy route}

This figure is included because the reader must see how Cross-domain decision question is constrained by Bounded answer plus residual risk. The highlighted review variable is option set, and the formal anchor is $V(K\mid E,R)$. The nearby prose names the proof, evidence row, table, replay check, or appendix that carries the claim.

\begin{figure}[!htbp]
\centering
\resizebox{0.96\textwidth}{!}{%
\begin{tikzpicture}[x=1cm,y=1cm,>=Latex,every node/.style={font=\small}]
  \definecolor{ocBlue}{HTML}{173B57}
\definecolor{ocGold}{HTML}{A77D2A}
\definecolor{ocGreen}{HTML}{4B7F52}
\definecolor{ocRed}{HTML}{8E3B32}
\definecolor{ocGray}{HTML}{ECEFF1}
  \draw[rounded corners=10pt,fill=ocGray!35,draw=ocBlue,line width=0.9pt] (-6.8,-3.4) rectangle (6.8,3.4);
  \node[font=\bfseries\large,ocBlue] at (0,3.0) {Strategy route};
  \node[draw,rounded corners=5pt,fill=blue!7,text width=0.24\textwidth,align=center,minimum height=1.05cm] (a) at (-4.35,1.0) {Cross-domain decision question};
  \node[draw,rounded corners=5pt,fill=green!8,text width=0.24\textwidth,align=center,minimum height=1.05cm] (b) at (4.35,1.0) {Bounded answer plus residual risk};
  \draw[->,line width=1.0pt,ocGreen] (a.east) -- node[above,fill=ocGray!35,inner sep=2pt,align=center] {decision compression} (b.west);
  \node[draw,rounded corners=5pt,fill=orange!10,text width=0.28\textwidth,align=center,minimum height=0.9cm] (r) at (-4.35,-1.45) {option set};
  \node[draw,rounded corners=5pt,fill=white,text width=0.28\textwidth,align=center,minimum height=0.9cm] (m) at (0,-1.45) {option set};
  \node[draw,rounded corners=5pt,fill=red!6,text width=0.28\textwidth,align=center,minimum height=0.9cm] (f) at (4.35,-1.45) {formal anchor: \( V(K\mid E,R) \)};
  \draw[->,dashed,ocGold] (r.north) -- (a.south);
  \draw[->,dashed,ocGold] (m.north) -- ($(a)!0.5!(b)$);
  \draw[->,dashed,ocGold] (f.north) -- (b.south);
\end{tikzpicture}}
\caption{Strategy route. This figure shows the reading relation from Cross-domain decision question to Bounded answer plus residual risk; the review variable is option set, the formal anchor is \( V(K\mid E,R) \), and the evidence link is the named proof, table, replay row, or appendix section cited near the figure.}
\label{fig:r012-inline-28e_oc133_inline_figures_reader_routes-4}
\end{figure}

\subsection{Evidence audit route}

This figure is included because the reader must see how Named evidence row is constrained by Checksum, source snapshot, and reviewer verdict. The highlighted review variable is hash/check, and the formal anchor is $H(\mathrm{source})=\mathrm{expected}$. The nearby prose names the proof, evidence row, table, replay check, or appendix that carries the claim.

\begin{figure}[!htbp]
\centering
\resizebox{0.96\textwidth}{!}{%
\begin{tikzpicture}[x=1cm,y=1cm,>=Latex,every node/.style={font=\small}]
  \definecolor{ocBlue}{HTML}{173B57}
\definecolor{ocGold}{HTML}{A77D2A}
\definecolor{ocGreen}{HTML}{4B7F52}
\definecolor{ocRed}{HTML}{8E3B32}
\definecolor{ocGray}{HTML}{ECEFF1}
  \draw[rounded corners=10pt,fill=ocGray!35,draw=ocBlue,line width=0.9pt] (-6.8,-3.4) rectangle (6.8,3.4);
  \node[font=\bfseries\large,ocBlue] at (0,3.0) {Evidence audit route};
  \node[draw,rounded corners=5pt,fill=blue!7,text width=0.24\textwidth,align=center,minimum height=1.05cm] (a) at (-4.35,1.0) {Named evidence row};
  \node[draw,rounded corners=5pt,fill=green!8,text width=0.24\textwidth,align=center,minimum height=1.05cm] (b) at (4.35,1.0) {Checksum, source snapshot, and reviewer verdict};
  \draw[->,line width=1.0pt,ocGreen] (a.east) -- node[above,fill=ocGray!35,inner sep=2pt,align=center] {verification} (b.west);
  \node[draw,rounded corners=5pt,fill=orange!10,text width=0.28\textwidth,align=center,minimum height=0.9cm] (r) at (-4.35,-1.45) {hash/check};
  \node[draw,rounded corners=5pt,fill=white,text width=0.28\textwidth,align=center,minimum height=0.9cm] (m) at (0,-1.45) {hash/check};
  \node[draw,rounded corners=5pt,fill=red!6,text width=0.28\textwidth,align=center,minimum height=0.9cm] (f) at (4.35,-1.45) {formal anchor: \( H(\mathrm{source})=\mathrm{expected} \)};
  \draw[->,dashed,ocGold] (r.north) -- (a.south);
  \draw[->,dashed,ocGold] (m.north) -- ($(a)!0.5!(b)$);
  \draw[->,dashed,ocGold] (f.north) -- (b.south);
\end{tikzpicture}}
\caption{Evidence audit route. This figure shows the reading relation from Named evidence row to Checksum, source snapshot, and reviewer verdict; the review variable is hash/check, the formal anchor is \( H(\mathrm{source})=\mathrm{expected} \), and the evidence link is the named proof, table, replay row, or appendix section cited near the figure.}
\label{fig:r012-inline-28e_oc133_inline_figures_reader_routes-5}
\end{figure}

\subsection{Venue projection route}

This figure is included because the reader must see how Release claim set is constrained by Journal-specific manuscript package without new claims. The highlighted review variable is checklist, and the formal anchor is $R_{\mathrm{release}}\rightarrow J_{\mathrm{venue}}$. The nearby prose names the proof, evidence row, table, replay check, or appendix that carries the claim.

\begin{figure}[!htbp]
\centering
\resizebox{0.96\textwidth}{!}{%
\begin{tikzpicture}[x=1cm,y=1cm,>=Latex,every node/.style={font=\small}]
  \definecolor{ocBlue}{HTML}{173B57}
\definecolor{ocGold}{HTML}{A77D2A}
\definecolor{ocGreen}{HTML}{4B7F52}
\definecolor{ocRed}{HTML}{8E3B32}
\definecolor{ocGray}{HTML}{ECEFF1}
  \draw[rounded corners=10pt,fill=ocGray!35,draw=ocBlue,line width=0.9pt] (-6.8,-3.4) rectangle (6.8,3.4);
  \node[font=\bfseries\large,ocBlue] at (0,3.0) {Venue projection route};
  \node[draw,rounded corners=5pt,fill=blue!7,text width=0.24\textwidth,align=center,minimum height=1.05cm] (a) at (-4.35,1.0) {Release claim set};
  \node[draw,rounded corners=5pt,fill=green!8,text width=0.24\textwidth,align=center,minimum height=1.05cm] (b) at (4.35,1.0) {Journal-specific manuscript package without new claims};
  \draw[->,line width=1.0pt,ocGreen] (a.east) -- node[above,fill=ocGray!35,inner sep=2pt,align=center] {format compliance} (b.west);
  \node[draw,rounded corners=5pt,fill=orange!10,text width=0.28\textwidth,align=center,minimum height=0.9cm] (r) at (-4.35,-1.45) {checklist};
  \node[draw,rounded corners=5pt,fill=white,text width=0.28\textwidth,align=center,minimum height=0.9cm] (m) at (0,-1.45) {checklist};
  \node[draw,rounded corners=5pt,fill=red!6,text width=0.28\textwidth,align=center,minimum height=0.9cm] (f) at (4.35,-1.45) {formal anchor: \( R_{\mathrm{release}}\rightarrow J_{\mathrm{venue}} \)};
  \draw[->,dashed,ocGold] (r.north) -- (a.south);
  \draw[->,dashed,ocGold] (m.north) -- ($(a)!0.5!(b)$);
  \draw[->,dashed,ocGold] (f.north) -- (b.south);
\end{tikzpicture}}
\caption{Venue projection route. This figure shows the reading relation from Release claim set to Journal-specific manuscript package without new claims; the review variable is checklist, the formal anchor is \( R_{\mathrm{release}}\rightarrow J_{\mathrm{venue}} \), and the evidence link is the named proof, table, replay row, or appendix section cited near the figure.}
\label{fig:r012-inline-28e_oc133_inline_figures_reader_routes-6}
\end{figure}

\clearpage
